\scanpage{68}
\mspage{64}
% source lines 1--2
\unsure{qui|que} implique aussitôt, pour $A,B\subset E$
\[
\bar\alpha(A)\mid\circ\mid\bar\alpha(B)\Longleftrightarrow A\mid\circ\mid B.
\]
% source lines 3--7
\unsure{Ces} conditions draconiennes donnent ce qu'il nous faut, mais peut-être même un peu \deleted{plus} trop --- \unsure{n'impliqueraient-elles} pas qu'on a
\diagram{\begin{tikzcd}[ampersand replacement=\&]
\alpha_*:\Sigma_E \arrow[r,"\sim"] \& \Sigma_{E'}\quad\text{bijectif ?}
\end{tikzcd}}
% source lines 8--14
\deleted{\illegible{2}}\inserted{D'autre} part, on aimerait peut-être \deleted{\unsure{mais|avoir}} avoir la formule $(*)$ p. 59, dans le cas où la formule $(**)$ p. 59 est fausse (comme ds l'exemple de tantôt). Donc je n'ai pas l'impression encore d'avoir vraiment bien compris la situation.

% source lines 15--17
Dans la formule $(*)$ p. 59
\[
\underbrace{\alpha_*\operatorname{supp}^{\unsure{\circ|\bullet}}_E(A)}_{\substack{\Vert\,\text{déf}\\\operatorname{supp}^{\circ}_{E'}\bar\alpha(\operatorname{supp}^{\circ}_E(A))}}
=\operatorname{supp}^{\circ}_{E'}(\bar\alpha(\unsure{A}))
\]
% source lines 18--21
l'inclusion $\supset$ est claire, puisque
\[
\operatorname{supp}^{\circ}_E(A)\supset\deleted{\operatorname{supp}^{\circ}_E}A\quad\text{d'où}\quad\deleted{\unsure{\alpha}}
\]
\[
\bar\alpha(\operatorname{supp}^{\circ}_E(A))\supset\bar\alpha(A)
\]
et il suffit de passer aux $\operatorname{supp}^{\circ}_{E'}$. Donc le p\textsuperscript{b} est
% archive sheet number: 67
